% Gaussian-posterior part of the exponential sample-schedule theorem.
% Requires theorem, lemma, proposition, corollary environments and amsmath,
% amssymb, amsthm. The spectral/random-Gram lemma is supplied separately.

\section{The actual Gaussian sign posterior}
Let $\mathcal X=\{-1,1\}^d$, $M=2^d$, and
$f(x)=\prod_{j=1}^d x_j$. All function norms and expectations on
$\mathcal X$ use its uniform probability measure. Let $G$ be a
nondegenerate centered Gaussian vector indexed by $\mathcal X$, with
covariance matrix $K$. The training inputs $X_1,\ldots,X_n$ are independent
and uniform. We condition on
\[
 \mathcal A_S=\{f(X_i)G(X_i)>0:1\le i\le n\}.
\]
The prediction at $x$ is its actual conditional correct-label probability
$q_S(x)=\mathbb P(f(x)G(x)>0\mid\mathcal A_S)$, and we write
\[
 F_\theta(S)=\mathbb P_x\{q_S(x)\le\theta\},\qquad
 \theta=\frac{13}{20},\qquad
 c=\sqrt{2\pi}\left(\frac1e-\frac7{20}\right)>0.
\]
The threshold $\theta<2/3$ supplies a strict buffer for passage from
limiting Gaussian posteriors to finite-width posterior probabilities.

Suppose the Walsh characters are eigenfunctions of the normalized
covariance operator, with degree-$j$ eigenvalue $\lambda_j$. Write
$V_r=\operatorname{span}\{\chi_U:|U|\le r\}$ and $D_r=\dim V_r$.
For $r<d$, let $Q$ be the covariance contributed by the complementary
Fourier modes, and put
\[
 \tau=Q(x,x)=\sum_{j>r}\binom dj\lambda_j,
 \qquad \ell=\max_{j>r}\lambda_j.
\]
The matrix eigenvalues of $Q$ are $M\lambda_j$; this normalization is
used throughout. Define $\|h\|_K^2=h^TK^{-1}h$, and write $u_+=\max(u,0)$ and $u_-=\max(-u,0)$.

\begin{lemma}[Mean norm and a confidence margin]\label{lem:schedule-mean}
Let $m(x)=\mathbb E[G(x)\mid\mathcal A_S]$, and let $s\le n$ be the
number of distinct sampled inputs. Then
\[
 \|m\|_K^2\le s,\qquad f(X_i)m(X_i)>0.
\]
Let $\sigma_x^2=\operatorname{Var}(G(x)\mid G_S)$, where $S$ here lists
the distinct inputs. Whenever $\sigma_x>0$ and $q_S(x)>\theta$,
\[
 f(x)m(x)>c\sigma_x.
\]
\end{lemma}
\begin{proof}
Whiten the Gaussian training vector as $G_S=K_{SS}^{1/2}\xi$.
The sign constraints form a homogeneous cone in $\xi$. Conditioning
leaves the squared Gaussian radius with mean $s$, so Jensen's inequality
and Gaussian regression give
\[
 m_S^TK_{SS}^{-1}m_S\le s,\qquad
 m=K_{\mathcal X,S}K_{SS}^{-1}m_S.
\]
The RKHS norm of this regression vector is the preceding quadratic
form. Its observed signed coordinates are strictly positive.

For the confidence claim, the scalar conditional variable
$Y=f(x)G(x)$ has a log-concave density: it is a marginal of a Gaussian
restricted to a convex set, by the marginal log-concavity theorem \cite{prekopa1973}. If $\mu=\mathbb E Y$ and $F_Y$ is its
cdf, log-concavity and Jensen's inequality imply
\[
 \log F_Y(\mu)\ge\mathbb E\log F_Y(Y)=-1.
\]
Consequently $\mathbb P(Y>\mu)\le1-e^{-1}<\theta$.
Conditioning first on the full training vector shows that the density
of $Y$ is a mixture of Gaussian densities with common variance
$\sigma_x^2$, and is therefore at most
$(\sqrt{2\pi}\sigma_x)^{-1}$. If $q_S(x)>\theta$, then $\mu>0$ and
\[
 \theta<\mathbb P(Y>0)
 \le 1-e^{-1}+\frac{\mu}{\sqrt{2\pi}\sigma_x}.
\]
This gives the stated margin.
\end{proof}

\begin{lemma}[Most residual variances remain large]\label{lem:schedule-var}
For every training sample,
\[
 \mathbb P_x\{\sigma_x^2<\tau/2\}\le \frac{2n\ell}{\tau}.
\]
Also, if $p$ is the degree-$r$ Fourier truncation of $m$ and $t=m-p$,
then $\|t\|_2^2\le n\ell$.
\end{lemma}
\begin{proof}
Write $G=G_{\le r}+G_{>r}$ as a sum of independent fields. Conditioning
additionally on the full low-degree field can only decrease conditional
Gaussian variance. The remaining covariance is
\[
 Q-Q_{\mathcal X,S}Q_{SS}^{\dagger}Q_{S,\mathcal X}.
\]
Its subtracted term is positive semidefinite, bounded above by $Q$,
and has rank at most $s\le n$. Its normalized trace is at most $n\ell$.
Markov's inequality proves the first assertion; it includes observed
points automatically. Fourier orthogonality and the mean norm bound
prove $\|t\|_2^2\le\ell\|m\|_K^2\le n\ell$.
\end{proof}

\section{Uniform sampling detects the negative part}
\begin{lemma}[Uniform ReLU sampling on the cube]\label{lem:schedule-relu}
For every fixed $\delta,\eta\in(0,1)$, there is a universal constant
$C$ such that, if
\[
 n\ge C\eta^{-2}D_r\log\!\left(\frac{D_r}{\delta\eta}\right),
\]
then with probability at least $1-\delta$, simultaneously for all
$p\in V_r$,
\[
 \left|\frac1n\sum_{i=1}^n(-f(X_i)p(X_i))_+
             -\mathbb E_x(-f(x)p(x))_+\right|
 \le\eta\mathbb E_x|p(x)|.
\]
In particular, when $r<d$ and $\eta=1/4$, the empirical negative
part is at least $\mathbb E|p|/4$.
\end{lemma}
\begin{proof}
Let $A$ have rows $f(x)(\chi_U(x))_{|U|\le r}$, one for each cube
point. Then $A^TA=MI$ and every row has squared norm $D_r$.
The $\ell_1$ Lewis weights of $A$ are exactly $D_r/M$: substitution
in their defining equation
\[
 w_x^2=A_x(A^T\operatorname{diag}(w^{-1})A)^{-1}A_x^T
\]
verifies this constant solution, and Lewis weights are unique.
The additive ReLU coreset theorem of Mai, Musco, and Rao
\cite[Theorem~5]{mai2021coresets} therefore samples precisely the uniform
distribution on the cube and gives the stated estimate. Its sample bound
is $O(\eta^{-2}D_r\log(D_r/(\delta\eta)))$; no additional logarithmic
factor is required. Finally $\mathbb E fp=0$, because $r<d$, so
$\mathbb E(-fp)_+=\mathbb E|p|/2$.
\end{proof}

\begin{proposition}[A deterministic reduction to a tail Gram bound]
\label{prop:schedule-deterministic}
Suppose a design obeys the conclusion of Lemma~\ref{lem:schedule-relu}
with $\eta=1/4$, and also
\begin{align}
 \|Q_{SS}\|_{\mathrm{op}}&\le B\tau,\label{eq:schedule-gram}\\
 \frac1n\#\{i:f(X_i)p(X_i)<0\}
       &\le\mathbb P_x\{f(x)p(x)\le0\}+u
       \quad\hbox{for all }p\in V_r.\label{eq:schedule-vc}
\end{align}
Here $B\ge1$, and the matrix in \eqref{eq:schedule-gram} lists the training sample
with multiplicity. Put $\zeta=n\ell/\tau$. If
\[
 u\le\frac{c^2}{2048B},\qquad
 \left(2+\frac8{c^2}\right)\zeta\le\frac{c^2}{2048B},
\]
then the actual Gaussian sign posterior satisfies
\[
 F_\theta(S)\ge\frac{c^2}{1024B}.
\]
\end{proposition}
\begin{proof}
Let $m=p+t$ be the low/high decomposition of its conditional mean,
and write $E_p=\mathbb P_x\{fp\le0\}$. Training consistency implies
\[
 (f(X_i)p(X_i))_-\le|t(X_i)|\mathbf1_{\{f(X_i)p(X_i)<0\}}.
\]
On the tail subspace $\|t\|_Q=\|t\|_K$. RKHS evaluation and $\|t\|_K^2\le n$ give
\[
 \frac1n\sum_i t(X_i)^2
 \le\frac{\|t\|_K^2}{n}\|Q_{SS}\|_{\mathrm{op}}\le B\tau.
\]
Combining the ReLU estimate, Cauchy--Schwarz, and
\eqref{eq:schedule-vc} yields
\begin{equation}\label{eq:schedule-l1upper}
 \mathbb E|p|\le4\sqrt{B\tau(E_p+u)}.
\end{equation}

Lemma~\ref{lem:schedule-var} and Markov's inequality show that, outside
sets of total measure at most
$a_0=(2+8/c^2)\zeta$, both
$\sigma_x\ge\sqrt{\tau/2}$ and
$|t(x)|\le c\sqrt\tau/(2\sqrt2)$ hold. At every remaining input
with confidence greater than $\theta$, Lemma~\ref{lem:schedule-mean}
gives
\[
 f(x)p(x)>\frac{c\sqrt\tau}{2\sqrt2}.
\]
Writing $e=F_\theta(S)+a_0$, we obtain $E_p\le e$ and
\[
 \frac{\mathbb E|p|}{\sqrt\tau}
 \ge\frac c{2\sqrt2}(1-e).
\]
If $e\le1/2$, comparison with \eqref{eq:schedule-l1upper} gives
$e+u\ge c^2/(512B)$, and hence
$F_\theta(S)\ge c^2/(1024B)$ under the stated assumptions.
If $e>1/2$, the same claimed lower bound follows directly from the
upper bound on $a_0$.
\end{proof}

\section{An exponentially large sample schedule}
The preceding proposition is deterministic. The remaining ingredients
are a standard VC estimate and the spectral/random-Gram estimates proved
in the preceding sections. We state their exact required form to make
the scope of the conclusion explicit.

\begin{theorem}[Failure along an exponential sample schedule]
\label{thm:exponential-schedule}
Fix the depth $L\ge2$. For the covariance
\[
 K(x,y)=H^{\circ(L-1)}\!\left(\frac{1+x\cdot y}{d+1}\right),
 \qquad H(t)=\frac2\pi\arcsin t,
\]
there are constants $A>0$, $\alpha\in(0,1/2)$, and $\varepsilon>0$
such that the following holds. Let $r_d$ be an odd integer with
$r_d/d\to\alpha$, and set
\[
 n_d=\left\lceil A D_{r_d}\log(D_{r_d}+1)\right\rceil.
\]
For all sufficiently large $d$, the Gaussian sign posterior trained on
$n_d$ independent uniform parity examples satisfies
\[
 \mathbb P_S\{F_{13/20}(S)\ge\varepsilon\}\ge\frac34.
\]
In particular $n_d=\exp((h(\alpha)+o(1))d)$, where
$h(\alpha)=-\alpha\log\alpha-(1-\alpha)\log(1-\alpha)>0$.
The same conclusion holds with $F_{13/20}(S)$ replaced by
$\liminf_{N\to\infty}F_N(S)$ in any finite-width model whose conditional
prediction probabilities converge to this Gaussian sign posterior.
This includes the negative-$\gamma$ scaling width limit.
\end{theorem}
\begin{proof}
Fix $B=2$ and $\varepsilon=c^2/(1024B)$. Set
$u=c^2/(2048B)$ and $\delta=1/20$. First choose $A$ large enough
for Lemma~\ref{lem:schedule-relu} with $\eta=1/4$ and this $\delta$.
The class $\{x: f(x)p(x)<0,\ p\in V_r\}$ has VC dimension at most
$D_r+1$. Thus the uniform VC theorem \cite{vapnik1971} gives
\eqref{eq:schedule-vc} with probability at least $1-\delta$ for all
large $d$, since $n_d/D_{r_d}\to\infty$.

The spectral estimates give, for odd $r$ and $r/d$ sufficiently small,
\[
 \frac{n_d\lambda_{r+1}}{\tau_r}
 \le C_L A\alpha^2\log(1/\alpha)+o(1).
\]
Corollary~\ref{cor:gram-schedule} gives
$\mathbb P_S\{\|Q_{SS}\|_{\mathrm{op}}>2\tau_r\}\le\delta$
for all large $d$ after $\alpha$ is made sufficiently small, depending
on $A,L,\delta$. Decrease $\alpha$ further so that the preceding spectral
ratio also meets Proposition~\ref{prop:schedule-deterministic}.
A union bound over the three sampling events proves a probability of
at least $1-3\delta>3/4$.

For the final assertion, fix $d$ and the sampled training set. At every
input with limiting correct-label probability at most $13/20$, the
finite-width correct-label probability is eventually below $2/3$.
The cube is finite, so this holds simultaneously at every such input,
which proves $\liminf_N F_N(S)\ge F_{13/20}(S)$.
\end{proof}

The theorem concerns the displayed sample schedule. It does not assert
failure at every smaller sample size, and no monotonicity of posterior
confidence in sample size is assumed. The later paper~\cite{gaussian-all-n} covers all intervening sample sizes
by an additional middle-band argument. The proof
also uses convexity of Gaussian sign conditioning; it does not establish
the analogous exponential conclusion at $\gamma=0$ or $0<\gamma<1$.
